% Flavor catalog process: C002
% family: charm
% owner_agent_id: PKA-C002
% writer_agent_id: null
% checker_agent_id: null
% opus_signoff_id: null
% Canonical metadata sidecar: C002.yaml
% Status is controlled by the YAML sidecar status_history, not this comment.

\section{\texorpdfstring{C002: CP violation in \(D^0-\bar D^0\) mixing}{C002: CP violation in D0-D0bar mixing}}

\subsection*{Process}
\textbf{Standard notation:} \(|q/p|\) and \(\phi_D=\arg(q/p)\) in
neutral-charm mixing.  This entry covers indirect CP violation in
\(D^0-\bar D^0\) oscillations; the CP-even mixing magnitudes
\(x_D\), \(y_D\), and \(\Delta m_D\) are cataloged separately in C001.

\subsection*{PDG / equivalent value(s)}
The canonical values are the HFLAV CKM25 all-CPV global-fit results, version
through 14 September 2025, accessed on 2026-05-16:
\[
  |q/p| = 0.983^{+0.015}_{-0.014},\qquad
  \phi_D = -1.51^{+1.03}_{-1.06}\ {\rm deg}.
\]
The corresponding 95\% C.L. intervals are \(|q/p|\in[0.955,1.012]\) and
\(\phi_D\in[-3.63,0.51]\) degrees.  HFLAV also reports that the no-indirect-CPV
point \((|q/p|,\phi)=(1,0)\) has \(\Delta\chi^2=2.16\), consistent with CP
conservation at \(0.95\sigma\).  These numbers are traced to
\texttt{flavor\_catalog/references/C002/hflav\_ckm25\_dmixing\_cpv\_fit.txt};
the PDG 2025 review excerpt in
\texttt{flavor\_catalog/references/C002/pdg2025\_dmixing\_review\_cpv.txt}
cross-checks the same central values.

\subsection*{Relevance to the RS / anarchic-flavor pipeline}
Anarchic warped models generate short-distance \(\Delta C=2\) amplitudes
through flavor-changing neutral-current exchange.  C002 is the CP-odd
counterpart of the D-mixing magnitude bound: it tests whether the complex
phase of a new \(c-u\) mixing amplitude is aligned well enough to avoid
visible indirect CP violation.  It is especially useful because the repo's
current D0 lane constrains \(|M_{12}^{D,{\rm NP}}|\), while C002 asks about the
phase-sensitive part of the same amplitude.

\subsection*{Post-2008 developments}
The baseline RS-flavor paper arXiv:0804.1954 predates the modern high-statistics
neutral-charm mixing program.  Since then, LHCb added CPV-sensitive
model-independent \(D^0\to K_S^0\pi^+\pi^-\) inputs, including the
post-Run-2 result summarized in
\texttt{lhcb2023\_arxiv2208\_06512.txt}.  Belle and Belle II have also supplied
new large-sample model-independent mixing information
(\texttt{belle\_belleii2025\_arxiv2410\_22961.txt}).  HFLAV now combines these
inputs in a global fit with explicit \(|q/p|\) and \(\phi_D\) parameters, and
the current result remains compatible with no indirect CP violation.

\subsection*{Constraint validity and model dependence}
This is a short-distance-sensitive but global-fit-dependent constraint.  The
experimental fit is precise enough to bound new CP phases, but the Standard
Model charm-mixing amplitude is long-distance dominated and difficult to
calculate.  A safe RS interpretation should therefore constrain the allowed
new-physics phase conservatively, without subtracting a precise SM prediction.

In down-aligned or kaon-protected RS variants, up-sector observables become
leading rather than secondary diagnostics.

\subsection*{Code coverage in this repo}
\textbf{PARTIAL.}  The legacy scan defines a D-mixing input at
\texttt{quarkConstraints/deltaf2.py:252} and evaluates D0 mixing through
\texttt{evaluate\_d0\_mixing} at \texttt{quarkConstraints/deltaf2.py:941}, but
that result uses \(|M_{12}^{\rm NP}|\) and returns only a magnitude ratio.
The scan exports \texttt{d\_mix\_ratio} at \texttt{quarkConstraints/scan.py:103}
and \texttt{quarkConstraints/scan.py:439}.  The modern policy labels D0 as
\texttt{conservative\_non\_cp\_mixing\_amplitude\_hadronic} at
\texttt{quarkConstraints/modern/phenomenology.py:46}.  A paper-mode helper does
retain complex \(M_{12}^{D0}\) and exposes \texttt{im\_M12\_D0\_NP\_GeV} at
\texttt{quarkConstraints/paper\_0710\_1869/eft\_deltaf2/observables.py:1403},
but there is no exported \(|q/p|\) or \(\phi_D\) evaluator.

\subsection*{Implementation difficulty}
\textbf{LOW.}  The existing \(\Delta F=2\) operator and RG machinery is the
right basis for C002, so no new operator class is required.  The missing step is
a CPV observable wrapper or likelihood policy that maps the complex D0 mixing
amplitude onto a conservative HFLAV \(|q/p|,\phi_D\) constraint.

\subsection*{Key references}
Process-local source keys before bibliography consolidation:
\texttt{hflav\_ckm25\_dmixing\_cpv\_fit},
\texttt{pdg2025\_dmixing\_review\_cpv},
\texttt{lhcb2023\_arxiv2208\_06512},
\texttt{belle\_belleii2025\_arxiv2410\_22961}, and
\texttt{cfw2008\_arxiv0804\_1954}.
